\documentclass[11pt,a4paper]{article}

\usepackage[utf8]{inputenc}
\usepackage[spanish]{babel}
\usepackage[a4paper,margin=1in]{geometry}
\usepackage{enumitem}
\usepackage{titlesec}
\usepackage{hyperref}

\titleformat{\section}
{\large\bfseries}
{}
{0em}
{}[\titlerule]

\begin{document}

\begin{center}
    {\LARGE \textbf{Juan Diego Cortes Barbosa}}\\[5pt]
    Junior Software Developer (Frontend)\\[5pt]
    \href{mailto:juandicortes07@gmail.com}{juandicortes07@gmail.com} \quad | \quad (+57) 3507427337
\end{center}

\vspace{0.5cm}

\section*{Resumen Profesional}

Tecnólogo en Análisis y Desarrollo de Software, recién graduado, con enfoque en el desarrollo frontend y gran interés en seguir fortaleciendo mis habilidades en diferentes áreas del desarrollo de software. Me destaco por crear interfaces modernas, responsivas y enfocadas en la experiencia del usuario, trabajando con tecnologías web actuales y buenas prácticas de desarrollo.

Cuento con bases en desarrollo frontend y backend, participando en proyectos académicos y personales donde he trabajado en la construcción de aplicaciones web, integración de funcionalidades y diseño de interfaces funcionales y escalables. Tengo facilidad para adaptarme a nuevos entornos, aprender tecnologías rápidamente y colaborar en equipo para cumplir objetivos de manera eficiente.

Soy una persona proactiva, comprometida y abierta a aprender lo que el proyecto o la empresa requiera, con interés en seguir creciendo profesionalmente y aportar soluciones de valor dentro de equipos de desarrollo.

\section*{Experiencia Profesional}

\textbf{Practicante de Desarrollo - Intergraficas S.A} \hfill 2025 -- 2026

\begin{itemize}[leftmargin=*]
    \item Participé en el desarrollo y mantenimiento de aplicaciones web utilizando Python con Django para la construcción de funcionalidades backend y lógica de negocio.
    
    \item Desarrollé interfaces frontend dinámicas y responsivas con React, enfocadas en mejorar la experiencia de usuario y la interacción con la plataforma.
    
    \item Trabajé con Strapi como CMS headless para la administración y gestión de contenido, integrando servicios y APIs entre frontend y backend.
    
    \item Colaboré en la integración de APIs y consumo de servicios para la comunicación entre diferentes módulos del sistema.
    
    \item Apoyé procesos de corrección de errores, validación de funcionalidades y mejora continua de la aplicación dentro del flujo de desarrollo.
    
    \item Participé en trabajo colaborativo con el equipo de desarrollo utilizando metodologías ágiles y buenas prácticas de programación.
\end{itemize}

\section*{Educación}

\textbf{Tecnólogo en Análisis y Desarrollo de Software}\\
Servicio Nacional de Aprendizaje (SENA)

\section*{Formación y Certificaciones}

\begin{itemize}[leftmargin=*]
    \item Lenguaje de programación Python -- MinTIC U. Distrital
    \item JavaScript Essentials -- Cisco
\end{itemize}

\end{document}